\section{The criteria on the normal forms are prefix conditions}\label{sec:combinatorics}

On a matching matrix, all three nullities entering the criteria are counts
of simple combinatorial objects, and the criteria collapse to conditions on
prefixes of the index set.

\subsection{The unsymmetric normal form}

Let $\Pi$ be a matching matrix with matching $\M$.  Relative to the cut
after index $k$, classify the matched pairs and define
\begin{align*}
  \iota_k &:= \#\{(i,j) \in \M : i \le k \text{ and } j \le k\}
  && \text{(pairs inside the leading block)},\\
  \chi^R_k &:= \#\{(i,j) \in \M : i \le k < j\}
  && \text{(pairs crossing the cut row-first)},\\
  \chi^C_k &:= \#\{(i,j) \in \M : j \le k < i\}
  && \text{(pairs crossing the cut column-first)},\\
  r^0_k &:= \#\{\text{zero rows} \le k\},
  \qquad
  c^0_k := \#\{\text{zero columns} \le k\}. \hspace{-12em}
\end{align*}

\begin{lemma}[nullities of a matching matrix]\label{lem:null-matching}
  For every $k$,
  \begin{equation}\label{eq:null-matching}
    \nl(\Pi[1:k,1:k]) = r^0_k + \chi^R_k = c^0_k + \chi^C_k,
    \qquad
    \nl(\Pi[:,1:k]) = c^0_k,
    \qquad
    \nl({\Pi[1:k,:]}^T) = r^0_k .
  \end{equation}
  In particular the \emph{bookkeeping identity}
  $r^0_k + \chi^R_k = c^0_k + \chi^C_k$ holds.
\end{lemma}

\begin{proof}
  The nonzero columns of the $n \by k$ block $\Pi[:,1:k]$ are the vectors
  $w_{ij} e_i$ over the matched pairs with $j \le k$; they occupy distinct
  rows and are therefore linearly independent, while the remaining $c^0_k$
  columns vanish.  Hence $\rk(\Pi[:,1:k]) = k - c^0_k$, giving the second
  formula, and the third follows by applying the same argument to $\Pi^T$,
  whose matching is the transpose of $\M$.

  The leading block $\Pi[1:k,1:k]$ is itself a matching matrix whose
  matched pairs are the $\iota_k$ pairs inside the block, so
  $\rk(\Pi[1:k,1:k]) = \iota_k$ (as in \Cref{prop:canonical}, its nonzero
  columns are independent).  Classify the $k$ row indices $\le k$: each is
  a zero row, or belongs to a pair with $j \le k$, or to a pair with
  $j > k$; the classes are disjoint and exhaust, so
  $k = r^0_k + \iota_k + \chi^R_k$ and
  $\nl(\Pi[1:k,1:k]) = k - \iota_k = r^0_k + \chi^R_k$.  Classifying the
  column indices instead gives $k = c^0_k + \iota_k + \chi^C_k$ and the
  alternative expression, whence the bookkeeping identity.
\end{proof}

\begin{lemma}[prefix form of the LU criterion]\label{lem:comb-lu}
  For a matching matrix $\Pi$, criterion \cref{eq:lu-crit} at index $k$ is
  equivalent to each of
  \[
    \text{(a) } \chi^R_k \le c^0_k,
    \qquad
    \text{(b) } \chi^C_k \le r^0_k,
    \qquad
    \text{(c) } \nu_k := \#\{(i,j) \in \M : \min(i,j) \le k\} \ \le\ k .
  \]
\end{lemma}

\begin{proof}
  Substituting \cref{eq:null-matching} into \cref{eq:lu-crit}: using the
  first expression for the leading nullity,
  $r^0_k + \chi^R_k \le c^0_k + r^0_k$, i.e.\ (a); using the second,
  $c^0_k + \chi^C_k \le c^0_k + r^0_k$, i.e.\ (b).  By the bookkeeping
  identity, (a) and (b) are equivalent (they have equal slack).  For (c):
  a pair has $\min(i,j) \le k$ iff it is counted by $\iota_k$, $\chi^R_k$,
  or $\chi^C_k$, so $\nu_k = \iota_k + \chi^R_k + \chi^C_k$; substituting
  $\iota_k = k - r^0_k - \chi^R_k$ from the proof of
  \Cref{lem:null-matching} gives
  $\nu_k \le k \iff \chi^C_k \le r^0_k$, which is (b).
\end{proof}

\subsection{The symmetric normal form}

Let $P$ be a symmetric matching matrix.  Define
\[
  \zeta_k := \#\{\text{isolated zero indices} \le k\},
  \qquad
  \chi_k := \#\{\text{pairs } (i,j) \text{ of } P : i \le k < j\} .
\]

\begin{lemma}[nullities of a symmetric matching matrix]\label{lem:null-sym-matching}
  For every $k$,
  \[
    \nl(P[1:k,1:k]) = \zeta_k + \chi_k,
    \qquad
    \nl(P[:,1:k]) = \zeta_k .
  \]
\end{lemma}

\begin{proof}
  The columns $t \le k$ of the $n \by k$ block $P[:,1:k]$ are: zero for
  isolated $t$; $\gamma_t e_t$ for singletons; for a pair $(i,j)$ of $P$,
  column $i$ of $P$ equals $\alpha e_j$ and column $j$ equals
  $\alpha e_i + \beta e_j$ (whichever of the two has index $\le k$ is
  included).  Distinct blocks occupy disjoint row sets, and within a pair
  with both $i, j \le k$ the two columns $\alpha e_j$ and
  $\alpha e_i + \beta e_j$ are independent (restricted to rows $\{i,j\}$
  they form the pair block, of determinant $-\alpha^2$).  Hence the nonzero
  columns of $P[:,1:k]$ are independent, exactly the isolated columns
  vanish, and $\nl(P[:,1:k]) = \zeta_k$.

  The leading block $P[1:k,1:k]$ is block diagonal (up to symmetric
  permutation) with blocks: zeros at isolated indices ($\zeta_k$ of them);
  $(\gamma_t)$ at singletons (rank $1$); full pair blocks when
  $j \le k$ (rank $2$); and, for each crossing pair $i \le k < j$, the
  single row/column $i$ whose entries within the block all vanish
  (its only nonzero, $\alpha$, sits in column $j > k$) --- contributing one
  null direction each.  Hence
  $\nl(P[1:k,1:k]) = \zeta_k + \chi_k$.
\end{proof}

\begin{lemma}[prefix form of the $LDL^T$ criterion]\label{lem:comb-sym}
  For a symmetric matching matrix $P$, criterion \cref{eq:ldlt-crit} at
  index $k$ is equivalent to
  \begin{equation}\label{eq:prefix-sym}
    \chi_k \ \le\ \zeta_k .
  \end{equation}
\end{lemma}

\begin{proof}
  Substitute \Cref{lem:null-sym-matching} into \cref{eq:ldlt-crit}:
  $\zeta_k + \chi_k \le 2 \zeta_k$.
\end{proof}

\begin{remark}\label{rem:comb-dictionary}
  Under the dictionary of \Cref{rem:compatibility} --- a pair $(i,j)$ of
  $P$ corresponds to the two matched pairs $(i,j)$ and $(j,i)$ of the
  unsymmetric normal form $\Pi$ of the same symmetric $A$, an isolated zero
  to a common zero row and column --- we have
  $\chi^R_k = \chi^C_k = \chi_k$ and $r^0_k = c^0_k = \zeta_k$, and
  conditions (a), (b) of \Cref{lem:comb-lu} both become
  \cref{eq:prefix-sym}.  The combinatorial obstruction is thus literally
  the same in the two problems.
\end{remark}
